%==============================================================================
% Chapter 4 — Bayesian ML & Probabilistic Modeling
% Curriculum: Phase 2, Chapter 4 of README.md
% Companion code: projects/04_bayesian_ml/
%
% Section skeleton only. Figures for this chapter go in theory/figures/.
%==============================================================================
\chapter{Bayesian ML \& Probabilistic Modeling}
\label{ch:bayesian-ml}

The models of \cref{ch:classical-ml} return a number.\projectref{projects/04\_bayesian\_ml/} The models here return a
distribution, and the difference matters whenever a prediction feeds a decision
with asymmetric costs. This chapter picks up \cref{eq:bayes} and confronts the
problem it created: the normalising integral in the denominator is intractable
for anything interesting, so the whole field is organised around ways to avoid
computing it.

Two escape routes get developed in full --- sample from the posterior without
normalising it (MCMC), or replace it with the closest tractable distribution
(variational inference). The evidence lower bound derived here returns in
\cref{ch:gen-ai} as the training objective for variational autoencoders, and
the Gaussian process machinery underpins the Bayesian hyperparameter
optimisation used throughout Phase~2.

%==============================================================================
\section{Core Bayesian Inference}
\label{sec:bayesian-inference}

\subsection{The Bayes Integral}

% TODO: Write out the full integral
% P(\theta|D) = P(D|\theta)P(\theta) / \int P(D|\theta)P(\theta) d\theta.

\subsection{Priors, Likelihoods and Posteriors}

\subsection{Conjugate Priors}

% TODO: Show how a Beta prior and a Binomial likelihood yield a Beta posterior.

%==============================================================================
\section{Markov Chain Monte Carlo}
\label{sec:mcmc}

\subsection{Detailed Balance}

\subsection{Metropolis--Hastings}

% TODO: Derive the acceptance probability.

\subsection{Hamiltonian Monte Carlo}

\subsection{Diagnostics: Convergence and Effective Sample Size}

%==============================================================================
\section{Variational Inference}
\label{sec:variational-inference}

\subsection{Posing Inference as Optimisation}

\subsection{Deriving the Evidence Lower Bound}

% TODO: Arrive at ELBO = E_q[\log P(X,Z) - \log q(Z)].

\subsection{Mean-Field Approximations}

%==============================================================================
\section{Gaussian Processes}
\label{sec:gaussian-processes}

\subsection{Distributions Over Functions}

\subsection{Covariance Functions and Kernels}

% TODO: RBF and Matern kernels.

\subsection{The Predictive Mean and Variance}

\subsection{Cost and Scaling}

%==============================================================================
\section{Bayesian Deep Learning}
\label{sec:bayesian-deep-learning}

\subsection{Bayesian Neural Networks}

\subsection{Monte Carlo Dropout as Approximate Inference}

% TODO: Prove how dropout approximates Bayesian inference.

\subsection{Calibration and Uncertainty Quality}
